\zlabel{triptych:concise:commentary:start}
\section{The Gospel received, practiced and shared}

\subsection{How the small beginning becomes fruitful}

Chrysostom's account of the mustard seed begins with the disciples' possible discouragement. They have heard about seed that fails and wheat endangered by weeds. Their difficulty concerns the people to whom they will preach: Christ has described resistance before telling them of growth. The mustard seed answers that difficulty by asking them to judge the beginning in relation to what Christ can bring from it. Their weakness is compatible with a work that exceeds their present strength.

Jerome's starting point concerns a different judgment. A hearer may prefer a polished philosophical discourse to preaching about a crucified God. What looks slight at first can nevertheless grow within a believing soul. Jerome allows the image of the sower to signify the Savior or the individual who receives the seed in his heart and nourishes it by faith. This second explanation makes the hearer's continued attention part of the growth. The branches signify teachings in which believing souls, or powers devoted to God, find rest. Understanding is among the fruits of the preaching.\footnote{Chrysostom, \textit{Homilies on Matthew}, 46.2; Jerome, \textit{Commentary on Matthew}, II, on Matthew 13:31--32, PL 26, cols.~89--91.}

Augustine explains why the grown mustard exceeds the other herbs. He identifies those herbs with the teachings accepted by different sects. The mustard's heat, or its reported power to expel poisons, signifies faith and its efficacy against those teachings. Chrysostom explains how preaching can reach many people; Augustine asks what makes faith effective against competing beliefs. Jerome's account of understanding concerns a person who must receive and live by that faith.\footnote{Augustine, \textit{Questions on the Gospels} (\textit{Quaestiones evangeliorum}), I.11, on Matthew 13:31--32.}

The difference becomes clearer in the leaven. Chrysostom reasons from the woman's action: she hides the leaven within the flour. Contact makes the change possible. Disciples who fear that being surrounded by opponents will destroy their work hear an assurance that Christ can make their presence fruitful there. The quantity of meal emphasizes the disparity between their apparent weakness and the effect of the Gospel. His explanation does not assign a separate spiritual identity to each measure.

Jerome instead explains what is affected throughout. In one interpretation the three measures signify spirit, soul and body; in another they signify reason, anger and desire. He identifies the woman with apostolic preaching or the Church gathered from the nations, and the leaven with knowledge and understanding of Scripture. The person receives this teaching through the Church. Jerome's account of inward change therefore includes both the giver of the teaching and the person who learns it. It can concern the Gospel's spread throughout the world as well as its growth within an individual.\footnote{Chrysostom, \textit{Homilies on Matthew}, 46.2, on Matthew 13:33; Jerome, \textit{Commentary on Matthew}, II, on 13:33, PL 26, cols.~91--92.}

Augustine identifies the woman with wisdom and the leaven with love that warms and rouses. One of his explanations of the three measures is the whole heart, soul and mind with which a person loves God. He also offers the three yields of fruit and the kinds of faithful people represented by Noah, Daniel and Job as alternatives. His leaven explains the affection that moves a person toward the good understood. This adds to Jerome's account of scriptural knowledge: recognizing virtue and being moved to pursue it are related parts of conversion.\footnote{Augustine, \textit{Questions on the Gospels}, I.12, on Matthew 13:33.}

These interpretations explain different features of the same comparison. Chrysostom emphasizes the small amount that changes the much greater amount around it. Jerome emphasizes the agreement produced within the person; Augustine explains love's power to stir that person to act. The Epistle permits a connection between them: the Thessalonians have themselves changed, and other people learn from their example. Their new allegiance becomes known through their lives. The explanations can complement one another without making contact, scriptural understanding and love the same account of the leaven.

At the literal level, Jesus describes growth and permeation. Chrysostom applies those actions to the disciples' mission; Jerome applies them to reception of the Gospel in the Church and the soul. Their spiritual explanations require a response from the hearer. At the end of the appointed passage, the parables invite further inquiry. Chrysostom limits the statement about Jesus speaking only in parables to this occasion; Jerome points to the disciples questioning him at home. The quoted Psalm 77(78):2 introduces mysteries within a history of God's dealings with his people. Learning to understand those dealings belongs to receiving the word that will be shared.\footnote{Matthew 13:34--36; Psalm 77(78):1--8; Chrysostom, \textit{Homilies on Matthew}, 47.1; Jerome, \textit{Commentary on Matthew}, II, on 13:34--35, PL 26, cols.~92--93.}

\subsection{The conduct through which faith is recognized}

Paul's thanksgiving concerns a community that has received power from God and has acted. Aquinas explains the power and fullness of the preaching in two ways. They can refer to its confirmation by miracles and the gift of the Holy Spirit. They can also refer to Paul's manner of preaching: a virtuous life, speech prompted by the Spirit and instruction sufficient for faith. The second explanation brings the preacher's conduct and teaching into the account of divine action. The first keeps the extraordinary gifts within it.\footnote{Aquinas, \textit{Commentary on First Thessalonians}, chapter 1, lecture 1 [87862], on 1:4--5.}

Chrysostom addresses a practical objection raised by the example of the apostles. His congregation might admire them while explaining their success entirely through miraculous powers. He replies that the possession of striking gifts does not establish the holiness of their possessor. The apostles' freedom from avarice and vainglory, and their mastery of the passions, matter to the lives they communicate. Listeners who cannot perform miracles can practice those virtues. His explanation of the leaven consequently leads to an examination of the people through whom the Gospel reaches others.\footnote{Chrysostom, \textit{Homilies on Matthew}, 46.3.}

Jerome examines the same practical problem through the faculties of the person. Understanding what is good does not by itself mean that anger and desire seek it. He assigns each a positive task. Reason is to attain prudence; anger is directed against vice; desire pursues virtue. Their agreement allows understanding to govern conduct. The Collect's request for words and deeds pleasing to God can be understood in this light: prayer asks for the integrity needed when knowledge of an obligation meets a contrary inclination.

This comparison makes the two interpretations more precise. In the first, cleansing and good conduct are needed because a community's life affects other people. In the second, those acts express the conversion of the person who performs them. A resentful word can expose a division between what someone understands about charity and what he wills toward his neighbor. The petition in the Collect reaches both the speaker and the person who hears him. God's help is sought for an action whose cause is inward and whose consequences are shared.

Aquinas further explains the three virtues Paul praises. He joins work and labor to faith: belief must act and continue acting under difficulty. He then describes abundant works of charity and the patience arising from hope. Chrysostom gives labor a particular connection with charity, explaining the cost of remaining attached to fellow believers under pressure. The difference in their grammatical exposition does not alter their agreement that the virtues can be recognized through conduct. Hope explains endurance because the good expected from Christ remains worth seeking when circumstances resist it.\footnote{Aquinas, \textit{Commentary on First Thessalonians}, 1, lecture 1 [87862], on 1:3; Chrysostom, \textit{Homilies on First Thessalonians}, homily 1, on 1:2--3; Jerome, \textit{Commentary on Matthew}, II, on 13:33, PL 26, col.~91.}

The passage from imitation to example gives that conduct a further consequence. Aquinas says Christians imitate teachers insofar as the teachers imitate Christ; their human failings are outside the reason for imitation. In this lesson the converts follow the apostolic example by receiving the word in tribulation. They then become an example to people who already believe. Their fidelity strengthens existing faith as well as making their conversion known. The moral sense of the mission interpretation includes responsible imitation and perseverance. The moral sense of Jerome's interpretation includes prudent judgment and educated desire. Both lead to actions by which someone else can be helped.

\subsection{Why the promise of peace leaves room for lament}

The Introit's promise and the Epistle's tribulation can initially seem difficult to hear together. Jeremiah's letter supplies a concrete setting for promised peace that requires patience. Its recipients are to build, plant and maintain households in exile. Jerome contrasts their expected restoration with false assurances of an immediate return. He also observes that God, though able to act without their assistance, calls them to prayer and wholehearted seeking. The recipients live toward the promised good through actions appropriate to their present situation.\footnote{Jeremiah 29:5--14; Jerome, \textit{Commentary on Jeremiah}, V, on 29:10--14, PL 24, cols.~893--894.}

The Gradual's full psalm poses a related difficulty from another direction. It recalls help already given and then laments the distress that has followed. Bellarmine explains the selected declaration of rescue as memory of God's former benefits. Augustine also interprets its accomplished deliverance prophetically: what God will certainly give can be spoken of as already done. Bellarmine's account begins from past rescue, Augustine's from the certainty of a future rescue. Both allow the worshipper to praise God while still needing deliverance. The psalm's insistence on continuing fidelity also prevents suffering from being equated with a particular personal offense.

Chrysostom explains the Spirit's joy within the Thessalonians' affliction by distinguishing the injury they suffer from their relation to Christ. What happens to them is grievous. The Spirit nevertheless sustains joy in belonging to the Lord for whose sake they endure. His comparison with the biblical children refreshed in the furnace locates that joy in God's action. Aquinas emphasizes their love for Christ as the reason they can bear hardship for him. Their explanations leave the pain real. This joy cannot be reduced to an outwardly cheerful manner that every distressed Christian must display.\footnote{Psalm 43(44):8--10,18--23; Bellarmine, \textit{Commentary on the Book of Psalms}, Psalm 43, explanations 7--9; Augustine, \textit{Expositions of the Psalms}, 43.6--8; Chrysostom, \textit{Homilies on First Thessalonians}, homily 1, on 1:6; Aquinas, \textit{Commentary on First Thessalonians}, 1, lecture 1 [87862], on 1:6.}

Augustine's treatment of the depths gives a different explanation of why a person who has heard God's promise continues to cry for help. He begins with Jonah: confinement cannot prevent prayer from reaching God. He then considers the person who appears prosperous but dismisses the need for mercy. Such a person may use his success as evidence that God is unconcerned with his sins. The one who cries has acknowledged both his need and a deliverer. Recognition of need is already a movement toward God.

Bellarmine distinguishes two truths required for penitence: the misery of sin and the mercy of God. Awareness of misery without mercy leads to despair; confidence without recognition of need leaves the person seeking no remedy. The full Psalm 129(130) supplies both the question of who could survive the reckoning of sins and the assurance of forgiveness. Augustine's explanation of Psalm 84(85) similarly interprets release from captivity through the forgiveness named immediately after the Introit's selected verse. Bellarmine distinguishes the liberty already received from its completion at the resurrection.\footnote{Augustine, \textit{Expositions of the Psalms}, 129.1--2 and 84.2--6; Bellarmine, \textit{Commentary on the Book of Psalms}, Psalm 129, explanations 3--5, and Psalm 84, explanation 1.}

The repeated prayer from the depths consequently belongs to both interpretations. Disciples who remain among opponents can seek help where escape is unavailable. People whose desires are being reformed can seek mercy before that reform is complete. The first explanation concerns perseverance in a difficult task; the second concerns the hopeful acknowledgment of sin. Neither makes faithful prayer dependent on already possessing the relief for which it asks.

\subsection{What the offering asks God to change}

At the Offertory, the petition from Psalm 129(130) returns after the Gospel has been heard. Its opening is repeated within the chant. The people bringing gifts still ask for mercy. Schuster explains their humility by observing that what they offer has itself come from God's goodness. Their act of giving therefore acknowledges a gift already received.\footnote{Schuster, \textit{The Sacramentary}, III, p.~188, on the Offertory \textit{De profundis}, ``Out of the depths.''}

Augustine's exposition proceeds beyond the opening verses appointed in the chants. He finds the sinner's reason for hope in the psalm's propitiation, which he explains through Christ's sacrifice. Forgiveness rests on what Christ has offered for those who cannot establish their innocence. This fuller psalm exposition illuminates the assembly's request at the offering. The Secret, on which Schuster comments directly, asks for four effects: cleansing, renewal, governance and protection. Their application to the people offering extends beyond the immediate removal of a fault.\footnote{Augustine, \textit{Expositions of the Psalms}, 129.3; Schuster, \textit{The Sacramentary}, I, p.~417.}

Read alongside Paul's thanksgiving, cleansing concerns wrongs that would contradict the faith a community professes. Renewal concerns its continued life in grace. Governance asks that its acts be rightly directed, and protection asks that the good received endure. Chrysostom's insistence on virtue in those who communicate the Gospel gives these petitions a practical urgency: their life influences the people among whom they remain.

Jerome's explanation of the person makes the requested governance more specific. Reason, anger and desire are active in the choices made after worship. Forgiveness does not remove the need to judge well, oppose vice and pursue virtue. Renewal must continue in the person's words and deeds. In this interpretation, the Secret's four requests describe what someone with divided purposes needs from God. Augustine's account of Christ's sacrifice explains why that person can ask confidently for a changed life.

The offering is followed by thanksgiving to the Father in the Trinity Preface. The living and true God whom the Thessalonians serve is confessed with the only-begotten Son and the Holy Spirit. The Preface names distinction of persons, unity of essence and equality of majesty, and joins the assembly's thanksgiving with the angels' praise. In the mission interpretation, the Gospel's spread is ordered toward this worship. In the interpretation of human conversion, understanding and desire are directed toward the God who has made himself known.

Jerome's reported Trinitarian explanation of the three measures needs its own qualification here. He calls its sense pious while denying that uncertain parabolic explanations can establish dogma. The three portions of flour and their mixing do not explain the divine persons or their unity. The Preface confesses what God has revealed. Jerome's caution permits the reader to receive his spiritual interpretations of the person without treating every possible correspondence in the image as equally demonstrative.\footnote{\textit{Missale Romanum} (1962), p.~293, no.~1082; Jerome, \textit{Commentary on Matthew}, II, on 13:33, PL 26, cols.~91--92.}

\subsection{Confidence in prayer and the good still awaited}

The Communion's assurance about asking and receiving raises a question about desire. A person may want something intensely while also needing that desire to change. Mark's context places faith in God beside the obligation to forgive. Bede explains faith working through charity, and directs those who have not attained extraordinary perfection to ask for pardon and eternal life. They are to forgive others who have sinned against them. His explanation concerns the connected Mark 11:22--25 passage, including the forgiveness command after the verse sung at Communion.\footnote{Mark 11:24--25; Bede, \textit{Commentary on Mark}, on 11:22--25, paragraphs beginning \textit{Et respondens Jesus} and \textit{Et cum stabitis ad orandum}.}

The obligation changes the petitioner's conduct toward another person. Someone asking mercy must be willing to extend it. Within the interpretation of inward conversion, this brings the spoken request and the intention of the will into agreement. Within the interpretation of public witness, it makes the community's life correspond to the mercy it professes. Forgiveness is part of the life sought in prayer; treating it as a technique for acquiring unrelated goods would miss Bede's point about charity.

The Postcommunion then asks that reception be followed by continued desire. Ordinary nourishment tends to end hunger, whereas the prayer asks the already fed to keep seeking the heavenly gifts by which they live. Schuster emphasizes this feature. The object of the desire is definite: the gifts already received. Their reception can deepen the recognition that life depends on them. In the first interpretation, this continuing dependence sustains a community's work and endurance. In the second, it describes desire being taught by grace.\footnote{Schuster, \textit{The Sacramentary}, I, p.~417, on the Postcommunion.}

Paul's final words identify the future for which the community waits: the Son raised from the dead, coming from heaven. Aquinas explains expectation through resurrection in conformity with Christ and deliverance from punishment through his freeing people from sin. The community's good reputation is a present fruit of conversion; its final hope is life with the risen Lord. Matthew's wider discourse likewise includes separation and judgment. The leaven's complete effect does not establish that every hearer will inevitably receive the Gospel.\footnote{1 Thessalonians 1:10; Aquinas, \textit{Commentary on First Thessalonians}, 1, lecture 1 [87862], on 1:10; Matthew 13:36--50.}

Augustine reaches this future through the full psalm behind the Alleluia and Offertory. He interprets the morning watch by Christ's resurrection: what has occurred in the Head gives the members reason to hope for their own bodily resurrection. Forgiveness already received and resurrection still awaited are related gifts. The flesh Christ raised is the humanity he assumed; those who belong to him may hope to share that life. Augustine consequently distinguishes Israel's final hope from making riches or bodily health the highest good. Temporal goods may be sought from God, while resurrection gives desire its further object.\footnote{Augustine, \textit{Expositions of the Psalms}, 129.4--8.}

This final, anagogical sense belongs to both interpretations with different emphases. The interpretation of witness under trial follows a community beyond its present influence to final deliverance with Christ. The interpretation of the whole person follows forgiveness and reformed desire toward complete redemption, including the body. Both retain a present task: the word is to be understood, practiced and communicated. Their confidence rests on the risen Son whose life has already been manifested and whose coming they await.
